\begin{textsample}{NEG Dim 3 – global\_south\_2024\_12 – Score -1.00 – global\_south\_2024\_12/119273830.txt}  \label{ex:f3_neg_069}
Hamas, Israel trade barbs over truce delay Hamas and Israel traded accusation on Wednesday over delays in finalising a ceasefire and hostage release deal, after both had reported progress in the latest round of Gaza truce talks.
Indirect negotiations mediated by Qatar, Egypt and the United States have taken place in Doha in recent days, rekindling hope for an agreement between the two warring parties that has so far proven elusive.
On Wednesday, both sides accused the other of throwing up roadblocks.
Palestinian militant group Hamas said in a statement that"the ceasefire and prisoner exchange negotiations are continuing in Doha under the mediation of Qatar and Egypt in a serious manner."But, it added, Israel"has set new conditions"which have"delayed reaching an agreement".
Hamas said the Israeli demands concern troop withdrawal, the terms of the proposed ceasefire and the potential release of prisoners held in Israeli jails, as well as"the return of displaced people"to their homes in Gaza.
It did not specify what the Israeli conditions were.
Israel swiftly refuted the accusations, saying it was Hamas that was creating"new obstacles"to an agreement."Hamas is once again lying, reneging on understandings already reached, and continuing to create new obstacles in the negotiations,"said a statement from the office of Prime \textbf{Minister} Benjamin Netanyahu.
Just days ago, both sides had reported progress.
Netanyahu on Monday told a parliament session that there had been"some progress"in the talks, and a day later his office said Israeli negotiators had returned from Qatar after"significant negotiations".
On Saturday, Hamas and two other Palestinian groups said in a joint statement that a ceasefire agreement was"closer than ever", but cautioned that it would only be possible if Israel did not impose new conditions.
In Israel, critics of Netanyahu including relatives of some of the dozens of hostages still in captivity in Gaza have accused the prime \textbf{minister} of stalling.
On Wednesday, a group of family members of Gaza hostages urged Netanyahu to secure a deal to bring their loved ones homes."It ' s time to bring them back, Netanyahu.
It ' s up to you,"said Sharon Sharabi, whose two brothers were abducted during the Hamas attack last year.
One of them has died in captivity, according to the Israeli military."It ' s about time, don't wait,"said Sharabi, reading out a statement at Tel Aviv ' s Hostages Square.
Efforts to strike a truce and hostage release deal in numerous rounds of indirect talks have repeatedly failed over key stumbling blocks.
Israel and Hamas have agreed just one truce in more than 14 months of war.

% matched lemmas: minister
\end{textsample}
